\section{Three-handle sphere closure}\label{sec:s-closure}

This section records the intended S argument.  Closed-manifold HJ
Theorem~5.3 is not used to conclude that $B$ is fixed.  Remove the final $4$-handle and call the
remaining handlebody $W_3$.  Its boundary is $S^3$.  Reversing the three
$3$-handles attaches three three-dimensional $1$-handles to $S^3$, and hence
\[
\partial W_2\cong\#^3(S^1\times S^2).
\]
The actual attaching sphere system has connected complement, since sphere
surgery on it gives the connected boundary of $W_3$.

\begin{lemma}[Kernel invariance under changing a complete sphere system]
\label{lem:Ssystem}
Let \(\mathcal S\) and \(\mathcal A\) be complete sphere systems in
\(\#^3(S^1\times S^2)\).  If they are related by isotopy, permutation and
sphere slides, then the kernels of the two total MWW three-handle maps out of
\(\mathcal S^2_0(W_2)\) agree.
\end{lemma}

\begin{proof}
An isotopy of an attaching sphere gives an isomorphic handle attachment.
A permutation merely reorders disjoint three-handles.  A sphere slide is
exactly a \(3\)--\(3\) handle slide.  Standard handle calculus gives a
diffeomorphism between the resulting three-handle cobordisms which is the
identity on their incoming boundary \(W_2\).  Naturality of the skein-lasagna
maps therefore gives \(q_{\mathcal A}=\Phi q_{\mathcal S}\) for an
isomorphism \(\Phi\) of the target modules.  Hence their kernels agree.
\end{proof}

Choose a standard complete system
\(\mathcal A=(A_1,A_2,A_3)\) in the Johnson replacement reversed
$1$-handle picture, already disjoint from the P0 reconstruction cube $B$.
Closed-manifold HJ Theorem~5.3 is not used to conclude that $B$ is fixed.
It is used only with Lemma~\ref{lem:Ssystem}: any other complete system,
including a going-up Kirby attaching system, is related by isotopy,
permutation and slides in the closed manifold, so the kernels agree.
The current version of \cite{HorvatJablonowski2025} (arXiv v3) contains
Lemma~5.5 (spotted-ball slide lemma) and Lemma~5.7 (relative uniqueness of
complete sphere systems); its Theorem~5.3 is proved through Lemma~5.7,
which in turn uses Lemma~5.5.  This paper invokes only Theorem~5.3 in the
closed manifold.  It does not invoke Lemma~5.7 in its relative form, which
would be the statement needed to keep the detector ball $B$ fixed during
the sphere slides; that relative statement is not used and the
corresponding fixed-$B$ conclusion is not claimed.
In the reversed $1$-handle picture, three belt spheres miss the P0 cube,
with dual loops pairing to the identity, $b=0$ endpoint foams, and
disjointness from the C1 model link and C2 supports; these are properties
of the model picture, and the corresponding statements for the actual
$\partial W_2$ are \Open\ (Proposition~\ref{thm:Sdischarge}).

MWW's intrinsic reformulation of a $3$-handle quotient uses the local module
action of $\mathcal S^2_0(S^2\times D^2)$.  At $N=2$, let $A_0$ denote the
essential sphere with one dot and $A_1$ the undotted sphere.  The quotient
relations are
\[
A_0=1,\qquad A_1=0.
\]
If $J_j$ is the equator of $A_j$, MWW's K\"unneth identification gives the
following exact description of the two hemisphere maps:
\[
\begin{array}{ccc}
\mathcal S^2_0(W_2;J_j)&\xrightarrow{\ (\Psi_{\Delta^+},\Psi_{\Delta^-})\ }&
\mathcal S^2_0(W_2)\oplus\mathcal S^2_0(W_2)\\
\big\downarrow\wr&&\big\downarrow{\ell\oplus\ell}\\
\mathcal S^2_0(W_2)\otimes\Q\{1,X\}
&\xrightarrow{\ (\ell\circ(1\otimes\epsilon),\,\ell\circ\mu_{A_j})\ }&
\Q\oplus\Q.
\end{array}
\tag{30}
\]
Here $\epsilon(X)=1$, $\epsilon(1)=0$, while
$\mu_{A_j}(v\otimes X)=vA_0$ and
$\mu_{A_j}(v\otimes1)=vA_1$.  Therefore S is equivalent to the two
whole-source identities
\[
\ell(vA_0)=\ell(v),\qquad \ell(vA_1)=0
\quad\text{for every }v\text{ and }j=1,2,3.
\tag{31}
\]
This identifies the formal maps with the coequalizer pair in MWW
Theorem~3.7, but not their endpoint evaluations for the actual essential
spheres.

\begin{proposition}[Fixed-detector sphere replacement]\label{thm:Sgeometry}
The total three-handle kernel may be computed using the standard belt-sphere
system of the Johnson replacement reversed $1$-handle picture, which misses
$B$, while the detector in $B$ remains fixed.  For this system, S reduces to
the six equations (31).
\end{proposition}

\begin{proof}
The reversed $1$-handle picture of the Johnson replacement places three
belt spheres missing $B$, with dual loops whose handle segments meet the
owner belt sphere once and whose chart returns miss every belt cube.
Closed-manifold HJ Theorem~5.3 relates any other complete
system to this one by isotopy in the closed manifold; Lemma~\ref{lem:Ssystem}
identifies the kernels.  The detector remains in $B$, which those belt
spheres miss.  The six equations (31) are then Lemma~\ref{lem:Sendpoint}.
All of this concerns the model reversed $1$-handle picture; the
corresponding constructions in the actual $\partial W_2$ are \Open\
(Proposition~\ref{thm:Sdischarge}).
\end{proof}

\subsection{The essential-sphere endpoint comparison}

\begin{lemma}[Endpoint factorization for a standard sphere]\label{lem:Sendpoint}
For every \(j\), every cable state and every source element \(v\), the
constant term of the actual endpoint image of
\(\mu_{A_j}(v\otimes a)\) is the old detector image of \(v\) tensored with
the rank-two TQFT map of the punctured standard sphere.  Consequently the
divided cubic row satisfies
\[
\ell(vA_0)=\ell(v),\qquad \ell(vA_1)=0.
\]
\end{lemma}

\begin{proof}
Make \(A_j\) transverse to the two-handle cocores, and let \(b_j\) be the
number of core disks.  MWW's proof of Theorem~3.10 represents the
\(\Delta^-\) hemisphere, before the core disks are restored, by the connected
genus-zero surface obtained from \(A_j\) after deleting those disks and the
\(\Delta^+\) disk.  Choose a generic movie after cutting the one-handles.
Because \(A_j\) and the detector sheet are disjoint, every critical point of
the movie belongs entirely to one of the two surfaces.  Mixed events are
therefore only invertible Reidemeister/pivotal moves, represented at the
endpoint by cable braids; there are no mixed births, saddles or deaths.

At \(q=1\), Lemma~\ref{lem:C2} identifies every mixed braid with the
symmetry of the tensor category.  Naturality of births, saddles and deaths
with respect to that symmetry moves all mixed braids to the ends of the
movie.  The same endpoint permutation and pivotal chart transport the
detector row and the surface map, so they cancel by simultaneous
conjugation.  Thus the actual constant map is
\[
\operatorname{Id}_{\rm old}\otimes\Delta^{b_j-1}\quad(b_j>0),
\]
and is \(\operatorname{Id}_{\rm old}\otimes\epsilon\) when \(b_j=0\).
The restored two-handle core disks give the \(b_j\) actual counits.  Hence
\[
\epsilon^{\otimes b_j}\Delta^{b_j-1}(1)=0,\qquad
\epsilon^{\otimes b_j}\Delta^{b_j-1}(X)=1
\tag{32}
\]
for \(b_j>0\), with the same equations interpreted as
\(\epsilon(1)=0,\epsilon(X)=1\) when \(b_j=0\).
The \(1,X\) basis is normalized by the \(\Delta^+\) cap in MWW
Example~3.8, and BHPW strict functoriality prevents an independent sign in
the \(\Delta^-\) movie.

The full mixed braid maps have the form \(P(I+O(h))\).  The permutation
\(P\) has just been cancelled, and pre- or postcomposing a row beginning in
order \(h^3\) with \(I+O(h)\) changes only order \(h^4\) and higher.
Therefore the constant calculation (32) is exactly the calculation seen by
the divided cubic row.  In the model picture the three belt spheres miss
the P0 cube and the C1 model $z$-circles, so $b_j=0$ and the foams are the
counit of MWW Example~3.8.  For the actual $\partial W_2$ the spheres
$A_j$, the numbers $b_j$, the intersections with the two-handle cocores and
the hemisphere movies are not constructed, so the lemma is established for
the model picture only.
\end{proof}

The endpoint exchange square is the $b=0$ counit of
Lemma~\ref{lem:Sendpoint}.
\[
\begin{array}{ccc}
\mathcal S^2_0(W_2;J_j)&\longrightarrow&
\mathcal S^2_0(W_2)\oplus\mathcal S^2_0(W_2)\\
\downarrow&&\downarrow\\
\mathcal S^2_0(W_2)\otimes\Q\{1,X\}&\longrightarrow&\Q\oplus\Q,
\end{array}
\]
whose lower maps are the actual counit row and actual \(A_j\)-action by the
movie decomposition in Lemma~\ref{lem:Sendpoint}.

\begin{proposition}[Status of S]\label{thm:Sdischarge}
Hypothesis~\ref{hyp:P2} is \Open\ for the Johnson replacement.
Established: Lemma~\ref{lem:Ssystem} (kernel invariance under isotopy,
permutation and sphere slides) and the algebraic form (30)--(32) on the
model picture.  Not established: an actual complete sphere system
$A_1,A_2,A_3$ in the actual $\partial W_2$ with the conditions of HJ
Theorem~5.3 verified (embedded, pairwise disjoint, connected complement,
spherical homology basis, disjoint from the detector ball, simultaneous
surgery giving $S^3$); the actual equators, hemispheres, intersections with
the two-handle cocores, restored core disks and Morse movies; and the
identities $D(vA_0)=D(v)$, $D(vA_1)=0$ for all source elements and all
relative-class summands computed from them.
\end{proposition}

\begin{proof}
Lemma~\ref{lem:Ssystem} is proved above.  The remaining items depend on
Hypotheses~\ref{hyp:P0} and~\ref{hyp:P1} and on constructions that are not
part of the current certificates.
\end{proof}
